% =====================================================================
\chapter{Soluciones de los ejercicios}
\label{ap:soluciones}
% =====================================================================

Cada solución se presenta en cinco partes: datos, planteamiento, desarrollo, resultado y comprobación. En los pasos intermedios se redondea a cuatro cifras significativas salvo que se indique otra cosa; los resultados exactos se dan como fracción cuando existe.

% ---------------------------------------------------------------------
\section{Tema 1}

\begin{solucion}{eje:1-contar}
\datos $K=\abs{\V}=5$; longitudes $T=0,1,2,3$.
\planteamiento por la \cref{prop:contar-secuencias}, hay $K^{T}$ secuencias de longitud $T$.
\pasos
\begin{enumerate}[label=\arabic*.]
  \item $K^{0}+K^{1}+K^{2}+K^{3}=1+5+25+125=156$.
  \item Con la fórmula: $(K^{4}-1)/(K-1)=(625-1)/4=624/4=156$.
  \item Longitud exactamente~3: $K^{3}=125$.
\end{enumerate}
\resultado \textbf{156} secuencias de longitud a lo sumo~3; \textbf{125} de longitud exactamente~3.
\comprobacion la secuencia vacía cuenta como una; $156-125=31=1+5+25$ son las de longitud menor que~3.
\end{solucion}

\begin{solucion}{eje:1-inyectiva}
\datos $\delta(\tau(s))=s$ para toda cadena $s\in\Sigma^{*}$.
\planteamiento para probar la inyectividad de $\tau$, se parte de $\tau(s)=\tau(s')$ y se aplica $\delta$.
\pasos
\begin{enumerate}[label=\arabic*.]
  \item Si $\tau(s)=\tau(s')$, aplicando $\delta$ a ambos lados: $s=\delta(\tau(s))=\delta(\tau(s'))=s'$. Luego $\tau$ es inyectiva.
  \item $\delta$ no tiene por qué ser inyectiva. Contraejemplo: si $\V$ contiene los tokens \texttt{a}, \texttt{b} y \texttt{ab}, las secuencias $(\texttt{a},\texttt{b})$ y $(\texttt{ab})$ son distintas y ambas se decodifican como la cadena «ab».
\end{enumerate}
\resultado \textbf{$\tau$ es inyectiva; $\delta$ en general no lo es}, aunque sí es sobreyectiva.
\comprobacion la sobreyectividad de $\delta$ se sigue de la misma condición: toda cadena $s$ es la imagen de $\tau(s)$.
\end{solucion}

\begin{solucion}{eje:1-granularidad}
\datos $10^{4}$ caracteres; $3{,}5$ caracteres por token en el tokenizador de subpalabras.
\planteamiento número de tokens $=$ número de caracteres $/$ caracteres por token.
\pasos
\begin{enumerate}[label=\arabic*.]
  \item Por caracteres: un token por carácter, $10^{4}$ tokens.
  \item Por subpalabras: $10^{4}/3{,}5=2857{,}1\approx 2857$ tokens.
  \item Con un contexto fijo de $C$ tokens, el tokenizador de subpalabras cubre $3{,}5\,C$ caracteres de texto, frente a $C$ caracteres por caracteres.
\end{enumerate}
\resultado \textbf{$10^{4}$ tokens frente a unos $2857$}. La tokenización por subpalabras permite que un contexto de longitud fija abarque unas $3{,}5$ veces más texto.
\comprobacion $2857\cdot 3{,}5=9999{,}5\approx 10^{4}$.
\end{solucion}

% ---------------------------------------------------------------------
\section{Tema 2}

\begin{solucion}{eje:2-casillas}
\datos $K=27$.
\planteamiento por la \cref{prop:parametros}, la tabla de $n$-gramas tiene $K^{n}$ entradas y $K^{n-1}(K-1)$ parámetros libres.
\pasos
\begin{enumerate}[label=\arabic*.]
  \item Bigramas: $K^{2}=27\cdot 27=729$ entradas; libres $27\cdot 26=702$.
  \item Trigramas: $K^{3}=729\cdot 27=19\,683$ entradas; libres $27^{2}\cdot 26=729\cdot 26=18\,954$.
\end{enumerate}
\resultado \textbf{729 entradas (702 libres) para bigramas y 19\,683 (18\,954 libres) para trigramas.}
\comprobacion cada token adicional de contexto multiplica la tabla por $K=27$: $729\cdot 27=19\,683$.
\end{solucion}

\begin{solucion}{eje:2-vf}
\planteamiento contrastar cada afirmación con la definición correspondiente.
\pasos
\begin{enumerate}[label=(\alph*)]
  \item \textbf{Falsa.} El \cref{teo:cadena} es una identidad que se deduce de la regla del producto. La aproximación aparece al imponer la hipótesis de Markov (\cref{def:ngramas}).
  \item \textbf{Verdadera.} Cada fila $i$ es la distribución $P(\cdot\mid i)$, luego $\sum_j P_{ij}=1$. En la \cref{fig:cadena}, la fila \texttt{a} da $3/5+1/5+1/5=1$.
\end{enumerate}
\resultado \textbf{(a) falsa; (b) verdadera.}
\end{solucion}

\begin{solucion}{eje:2-sol-matriz}
\datos corpus «sol», «sal», «los»; $\V=\{\fin,\texttt{a},\texttt{l},\texttt{o},\texttt{s}\}$.
\planteamiento añadir \fin\ al principio y al final de cada palabra, extraer los bigramas, contar y dividir cada fila por su total (\cref{teo:emv}).
\pasos
\begin{enumerate}[label=\arabic*.]
  \item Secuencias: \fin sol\fin, \fin sal\fin, \fin los\fin.
  \item Bigramas (12 en total):
  \[\begin{aligned}
  &\text{sol: } (\fin,s),\,(s,o),\,(o,l),\,(l,\fin);\qquad \text{sal: } (\fin,s),\,(s,a),\,(a,l),\,(l,\fin);\\
  &\text{los: } (\fin,l),\,(l,o),\,(o,s),\,(s,\fin).
  \end{aligned}\]
  \item Matriz de conteos y matriz de transición:
\end{enumerate}
\begin{center}\small
\begin{tabular}{@{}c|ccccc|c@{}}
\toprule
$N$ & \fin & a & l & o & s & total \\ \midrule
\fin & 0 & 0 & 1 & 0 & 2 & 3 \\
a & 0 & 0 & 1 & 0 & 0 & 1 \\
l & 2 & 0 & 0 & 1 & 0 & 3 \\
o & 0 & 0 & 1 & 0 & 1 & 2 \\
s & 1 & 1 & 0 & 1 & 0 & 3 \\ \bottomrule
\end{tabular}
\qquad
\begin{tabular}{@{}c|ccccc@{}}
\toprule
$P$ & \fin & a & l & o & s \\ \midrule
\fin & 0 & 0 & $1/3$ & 0 & $2/3$ \\
a & 0 & 0 & $1$ & 0 & 0 \\
l & $2/3$ & 0 & 0 & $1/3$ & 0 \\
o & 0 & 0 & $1/2$ & 0 & $1/2$ \\
s & $1/3$ & $1/3$ & 0 & $1/3$ & 0 \\ \bottomrule
\end{tabular}
\end{center}
\resultado \textbf{las dos matrices anteriores}; por ejemplo, $P(\texttt{s}\mid\fin)=2/3$ y $P(\texttt{l}\mid\texttt{a})=1$.
\comprobacion los totales suman $3+1+3+2+3=12$ bigramas y cada fila de $P$ suma~1.
\end{solucion}

\begin{solucion}{eje:2-sol-prob}
\datos la matriz $P$ del \cref{eje:2-sol-matriz}.
\planteamiento por el \cref{cor:cadena-secuencias}, la probabilidad de una palabra es el producto de sus transiciones.
\pasos
\begin{enumerate}[label=\arabic*.]
  \item $P(\text{sol})=P(\texttt{s}\mid\fin)P(\texttt{o}\mid\texttt{s})P(\texttt{l}\mid\texttt{o})P(\fin\mid\texttt{l})=\frac23\cdot\frac13\cdot\frac12\cdot\frac23=\frac{4}{54}=\frac{2}{27}$.
  \item $P(\text{sal})=\frac23\cdot\frac13\cdot 1\cdot\frac23=\frac{4}{27}$.
  \item $P(\text{los})=P(\texttt{l}\mid\fin)P(\texttt{o}\mid\texttt{l})P(\texttt{s}\mid\texttt{o})P(\fin\mid\texttt{s})=\frac13\cdot\frac13\cdot\frac12\cdot\frac13=\frac{1}{54}$.
  \item $P(\text{sos})=\frac23\cdot\frac13\cdot\frac12\cdot\frac13=\frac{2}{54}=\frac{1}{27}$.
  \item Cociente: $P(\text{sal})/P(\text{sol})=(4/27)/(2/27)=2$. Las dos palabras solo difieren en el tramo central: $\texttt{s}\to\texttt{a}\to\texttt{l}$ vale $\frac13\cdot 1=\frac13$ y $\texttt{s}\to\texttt{o}\to\texttt{l}$ vale $\frac13\cdot\frac12=\frac16$.
\end{enumerate}
\resultado \textbf{$P(\text{sol})=2/27\approx 0{,}07407$; $P(\text{sal})=4/27\approx 0{,}1481$; $P(\text{los})=1/54\approx 0{,}01852$; $P(\text{sos})=1/27\approx 0{,}03704$.}
\comprobacion «sos» no está en el corpus y tiene probabilidad positiva; la suma de las cuatro es menor que~1, como corresponde a cuatro palabras entre infinitas.
\end{solucion}

\begin{solucion}{eje:2-longitud}
\datos $P(\text{a})=0{,}40$, $P(\text{ana})=0{,}08$, $P(\text{luna})=0{,}10$ (\cref{ej:bigramas}).
\planteamiento cada transición multiplica por un factor $\le 1$; para comparar longitudes distintas se usa la log-probabilidad media por transición, $\frac1n\ln P$.
\pasos
\begin{enumerate}[label=\arabic*.]
  \item «a» tiene 2 transiciones: $\ln 0{,}40/2=-0{,}9163/2=-0{,}4581$.
  \item «ana» tiene 4: $\ln 0{,}08/4=-2{,}526/4=-0{,}6314$.
  \item «luna» tiene 5: $\ln 0{,}10/5=-2{,}303/5=-0{,}4605$.
\end{enumerate}
\resultado \textbf{«a» es la más probable solo por ser la más corta}. Por transición, «a» ($-0{,}4581$) y «luna» ($-0{,}4605$) prácticamente empatan y «ana» queda por detrás. La medida propuesta es la pérdida del \cref{tema:5} cambiada de signo.
\comprobacion en perplejidad por transición, $e^{0{,}4581}\approx 1{,}581$ y $e^{0{,}4605}\approx 1{,}585$.
\end{solucion}

\begin{solucion}{eje:2-sin-fin}
\datos cuatro frases sin símbolo de fin; vocabulario $\{\texttt{a},\texttt{b}\}$ más $\langle s\rangle$.
\planteamiento contar las transiciones desde cada token que tiene sucesor y normalizar; después multiplicar.
\pasos
\begin{enumerate}[label=\arabic*.]
  \item Desde $\langle s\rangle$: a (frases 1 y 4) y b (frases 2 y 3); $P(\texttt{a}\mid\langle s\rangle)=P(\texttt{b}\mid\langle s\rangle)=1/2$.
  \item Desde \texttt{a} (solo cuentan las \texttt{a} seguidas de otro token, frases 1 y 4): $\texttt{a}\to\texttt{b}$ dos veces; $P(\texttt{b}\mid\texttt{a})=1$, $P(\texttt{a}\mid\texttt{a})=0$.
  \item Desde \texttt{b} (frases 2 y 3): $\texttt{b}\to\texttt{b}$ y $\texttt{b}\to\texttt{a}$; $P(\texttt{b}\mid\texttt{b})=P(\texttt{a}\mid\texttt{b})=1/2$.
  \item (b) $P(\texttt{aa})=\frac12\cdot 0=0$; $P(\texttt{ab})=\frac12\cdot 1=\frac12$; $P(\texttt{ba})=\frac12\cdot\frac12=\frac14$; $P(\texttt{bb})=\frac14$. Suma: $0+\frac12+\frac14+\frac14=1$.
  \item (c) Es el caso $L=2$ de la \cref{prop:longitud-fija}: cada longitud suma~1, de modo que la suma sobre todas las longitudes diverge.
\end{enumerate}
\resultado \textbf{$P(\texttt{aa})=0$, $P(\texttt{ab})=1/2$, $P(\texttt{ba})=P(\texttt{bb})=1/4$}; el modelo sin símbolo de fin no es un modelo de lenguaje en el sentido de la \cref{def:ml}.
\comprobacion las ocho frases de longitud~3 también suman~1 (por ejemplo, $P(\texttt{abb})=\frac12\cdot 1\cdot\frac12=\frac14$).
\end{solucion}

% ---------------------------------------------------------------------
\section{Tema 3}

\begin{solucion}{eje:3-laplace}
\datos conteos del \cref{eje:2-sol-matriz}; $\alpha=1$; $K=5$.
\planteamiento $\widehat P_1(v\mid u)=(N(u,v)+1)/(N(u)+5)$ (\cref{def:laplace}).
\pasos
\begin{enumerate}[label=\arabic*.]
  \item $P(\fin\mid\fin)$: $N(\fin,\fin)=0$ y $N(\fin)=3$; $P=(0+1)/(3+5)=1/8$.
  \item $P(\texttt{l}\mid\texttt{a})$: $N(\texttt{a},\texttt{l})=1$ y $N(\texttt{a})=1$; $P=(1+1)/(1+5)=2/6=1/3$.
\end{enumerate}
\resultado \textbf{$P(\fin\mid\fin)=1/8=0{,}125$ (antes~0) y $P(\texttt{l}\mid\texttt{a})=1/3\approx 0{,}3333$ (antes~1).}
\comprobacion la fila \texttt{a} suavizada es $(1,1,2,1,1)/6$, que suma $6/6=1$.
\end{solucion}

\begin{solucion}{eje:3-u1}
\datos cuatro frases; $\alpha=1$.
\planteamiento $P(\text{Ana}\mid\text{soy})=(N(\text{soy},\text{Ana})+1)/(N(\text{soy})+K)$.
\pasos
\begin{enumerate}[label=\arabic*.]
  \item Vocabulario: $\{\langle s\rangle,\text{yo},\text{soy},\text{Ana},\langle/s\rangle,\text{no},\text{quiero},\text{sopa},\text{ni}\}$, luego $K=9$.
  \item «soy» aparece en las frases 1, 2 y 3, siempre con sucesor: Ana (frases 1 y 3) o $\langle/s\rangle$ (frase 2). Así, $N(\text{soy})=3$ y $N(\text{soy},\text{Ana})=2$.
  \item $P=(2+1)/(3+9)=3/12=1/4$.
\end{enumerate}
\resultado \textbf{$P(\text{Ana}\mid\text{soy})=1/4=0{,}25$} (sin suavizar, $2/3\approx 0{,}6667$).
\comprobacion la fila «soy» suavizada asigna $3/12$ a Ana, $2/12$ a $\langle/s\rangle$ y $1/12$ a cada uno de los otros 7 tokens: $3+2+7=12$.
\end{solucion}

\begin{solucion}{eje:3-interp}
\datos $N(\text{soy},\text{Ana})=2$, $N(\text{soy})=3$, $N(\text{Ana})=4$; total de tokens $5+5+5+8=23$; $\lambda_1=\lambda_2=1/2$.
\planteamiento $P_{\mathrm{int}}=\tfrac12\widehat P_2(\text{Ana}\mid\text{soy})+\tfrac12\widehat P_1(\text{Ana})$.
\pasos
\begin{enumerate}[label=\arabic*.]
  \item Bigrama: $\widehat P_2=2/3$.
  \item Unigrama: «Ana» aparece una vez en cada frase, $\widehat P_1=4/23$.
  \item $P_{\mathrm{int}}=\tfrac12\cdot\tfrac23+\tfrac12\cdot\tfrac{4}{23}=\tfrac13+\tfrac{2}{23}=\tfrac{23}{69}+\tfrac{6}{69}=\tfrac{29}{69}$.
\end{enumerate}
\resultado \textbf{$P(\text{Ana}\mid\text{soy})=29/69\approx 0{,}4203$.}
\comprobacion el valor está entre el unigrama ($0{,}1739$) y el bigrama ($0{,}6667$), como corresponde a una media ponderada.
\end{solucion}

\begin{solucion}{eje:3-limite}
\datos $u$ fijo con $n=N(u)>0$; $\widehat P_\alpha(v\mid u)=(N(u,v)+\alpha)/(n+\alpha K)$.
\planteamiento límites de un cociente de funciones continuas en $\alpha$.
\pasos
\begin{enumerate}[label=\arabic*.]
  \item $\alpha\to 0^{+}$: numerador $\to N(u,v)$ y denominador $\to n>0$; el cociente tiende a $N(u,v)/n$, que es el EMV del \cref{teo:emv}.
  \item $\alpha\to\infty$: dividiendo numerador y denominador por $\alpha$, $\widehat P_\alpha=(N(u,v)/\alpha+1)/(n/\alpha+K)\to 1/K$.
  \item Interpretación: $\widehat P_\alpha$ es la frecuencia relativa que se obtendría si, además de los datos reales, se hubiera observado $\alpha$ veces cada uno de los $K$ tokens tras $u$. Es la media a posteriori con prior de Dirichlet (\cref{prop:laplace}).
\end{enumerate}
\resultado \textbf{$\lim_{\alpha\to 0^+}\widehat P_\alpha=N(u,v)/N(u)$ y $\lim_{\alpha\to\infty}\widehat P_\alpha=1/K$}. Un $\alpha$ grande da más peso a la creencia previa uniforme que a los datos.
\comprobacion con la fila \texttt{a} del \cref{eje:3-laplace} y $\alpha=100$: $P(\texttt{l}\mid\texttt{a})=101/501\approx 0{,}2016$, cerca de $1/5$.
\end{solucion}

% ---------------------------------------------------------------------
\section{Tema 4}

\begin{solucion}{eje:4-orden}
\planteamiento aplicar la propiedad (iii) de la \cref{prop:softmax}.
\pasos
\begin{enumerate}[label=\arabic*.]
  \item Para todo $T>0$, $s\mapsto e^{s/T}$ es estrictamente creciente.
  \item Si $z_i>z_j$, entonces $e^{z_i/T}>e^{z_j/T}$ y, al dividir por el mismo denominador, $p_i>p_j$.
\end{enumerate}
\resultado \textbf{No.} La temperatura modifica cuánto destaca el máximo, no cuál es.
\comprobacion en el \cref{tab:temperatura}, $p_1$ es la mayor con $T=0{,}5$, $1$ y $2$.
\end{solucion}

\begin{solucion}{eje:4-softmax}
\datos $z=(0,\ln 2,\ln 5)$; $T=1$ y $T=2$.
\planteamiento \cref{def:softmax}, usando $e^{\ln a}=a$ y $e^{(\ln a)/2}=\sqrt a$.
\pasos
\begin{enumerate}[label=\arabic*.]
  \item $T=1$: $e^{z}=(1,2,5)$, suma $8$; $\softmax(z)=(1/8,\,2/8,\,5/8)=(0{,}125;\,0{,}25;\,0{,}625)$.
  \item $T=2$: $e^{z/2}=(1,\sqrt2,\sqrt5)=(1;\,1{,}414;\,2{,}236)$, suma $4{,}650$.
  \item Cocientes: $1/4{,}650=0{,}2151$; $1{,}414/4{,}650=0{,}3041$; $2{,}236/4{,}650=0{,}4809$.
  \item $T\to 0^{+}$: el máximo $z_3=\ln 5$ es único; por la propiedad (iv), el límite es $(0,0,1)$.
\end{enumerate}
\resultado \textbf{$T=1$: $(0{,}125;\,0{,}25;\,0{,}625)$. $T=2$: $(0{,}2151;\,0{,}3041;\,0{,}4809)$. $T\to 0^{+}$: $(0,0,1)$.}
\comprobacion la segunda fila suma $0{,}2151+0{,}3041+0{,}4809=1{,}000$; al pasar de $T=1$ a $T=2$ el máximo baja de $0{,}625$ a $0{,}4809$.
\end{solucion}

\begin{solucion}{eje:4-muestreo}
\datos $p=(0{,}1;\,0{,}5;\,0{,}4)$; $u\in\{0{,}05;\,0{,}35;\,0{,}6;\,0{,}95\}$.
\planteamiento $X=\min\{k:u<F_k\}$ (\cref{prop:muestreo}).
\pasos
\begin{enumerate}[label=\arabic*.]
  \item $F_1=0{,}1$; $F_2=0{,}6$; $F_3=1{,}0$. Intervalos: $x\leftrightarrow[0;0{,}1)$, $y\leftrightarrow[0{,}1;0{,}6)$, $z\leftrightarrow[0{,}6;1)$.
  \item $u=0{,}05<0{,}1$: \textbf{x}. $u=0{,}35$: $0{,}1\le 0{,}35<0{,}6$, \textbf{y}.
  \item $u=0{,}6$: $0{,}6$ no es menor que $F_2=0{,}6$, pero sí que $F_3$: \textbf{z}. $u=0{,}95$: \textbf{z}.
\end{enumerate}
\resultado \textbf{x, y, z, z.}
\comprobacion la longitud de cada intervalo coincide con su probabilidad: $0{,}1$, $0{,}5$ y $0{,}4$.
\end{solucion}

\begin{solucion}{eje:4-potencia}
\datos $Q(i)\propto P(i)^{\beta}$; $P=(0{,}5;\,0{,}25;\,0{,}25)$.
\planteamiento escribir $P(i)^{\beta}=e^{\beta\ln P(i)}=e^{z_i/T}$ con $z_i=\ln P(i)$ y $T=1/\beta$.
\pasos
\begin{enumerate}[label=\arabic*.]
  \item (a) Duplicar la temperatura ($T=2$) es $\beta=1/2$; reducirla a la mitad ($T=1/2$) es $\beta=2$.
  \item (b) $\beta=2$: $P^{2}=(0{,}25;\,0{,}0625;\,0{,}0625)$, suma $0{,}375$; $Q=(2/3;\,1/6;\,1/6)$.
  \item $\beta=1/2$: $\sqrt{P}=(0{,}7071;\,0{,}5;\,0{,}5)$, suma $1{,}707$; $Q=(0{,}4142;\,0{,}2929;\,0{,}2929)$.
\end{enumerate}
\resultado \textbf{$\beta=1/T$}; con $\beta=2$, $Q=(0{,}6667;\,0{,}1667;\,0{,}1667)$; con $\beta=1/2$, $Q=(0{,}4142;\,0{,}2929;\,0{,}2929)$.
\comprobacion ambas distribuciones suman~1 y conservan el orden de $P$.
\end{solucion}

\begin{solucion}{eje:4-limites}
\datos $z\in\R^{K}$, $T>0$, $c\in\R$.
\planteamiento manipular el cociente y tomar límites término a término.
\pasos
\begin{enumerate}[label=\arabic*.]
  \item (ii) $e^{(z_i+c)/T}=e^{z_i/T}e^{c/T}$ por $e^{a+b}=e^{a}e^{b}$. El factor $e^{c/T}>0$ aparece en el numerador y en cada sumando del denominador; al sacarlo como factor común se cancela.
  \item (iv) Sea $z_k$ el máximo único. Con $c=-z_k$: $p_i=e^{(z_i-z_k)/T}/\sum_j e^{(z_j-z_k)/T}$. Para $j=k$ el sumando vale~1; para $j\ne k$, $z_j-z_k<0$, luego $(z_j-z_k)/T\to-\infty$ y $e^{(z_j-z_k)/T}\to 0$. El denominador tiende a~1, $p_k\to 1$ y $p_i\to 0$ si $i\ne k$.
  \item (v) Para $z$ fijo, $z_j/T\to 0$ y, por continuidad, $e^{z_j/T}\to 1$; luego $p_i\to 1/K$.
  \item Contraejemplo para (iv): $z=(1,1,0)$. Con $c=-1$, $p=(1,1,e^{-1/T})/(2+e^{-1/T})\to(1/2,1/2,0)$, que no es un vector canónico.
\end{enumerate}
\resultado \textbf{Las tres propiedades quedan demostradas; sin máximo único, el límite reparte la masa entre los máximos.} $\blacksquare$
\comprobacion con $z=(2,1,0)$ y $T=0{,}1$, $\softmax_T(z)\approx(0{,}99995;\,4{,}54\cdot 10^{-5};\,2\cdot 10^{-9})$; con $T=100$, $\approx(0{,}3367;\,0{,}3333;\,0{,}3300)$.
\end{solucion}

\begin{solucion}{eje:4-lse}
\datos $g(x)=\ln\sum_i e^{x_i}$; $x=(1000;\,1001;\,1002)$.
\planteamiento regla de la cadena para (a); factor común para (b); (b) con $b=-\max_i x_i$ para (c).
\pasos
\begin{enumerate}[label=\arabic*.]
  \item (a) Con $\frac{d}{du}\ln u=\frac1u$: $\partial g/\partial x_i=e^{x_i}/\sum_j e^{x_j}=\softmax(x)_i$.
  \item (b) $\sum_i e^{x_i+b}=e^{b}\sum_i e^{x_i}$; tomando logaritmos, $g(x+b\mathbf 1)=b+g(x)$.
  \item (c) Con $m=1002$: $g(x)=1002+\ln(e^{-2}+e^{-1}+e^{0})$.
  \item $e^{-2}=0{,}1353$; $e^{-1}=0{,}3679$; suma $0{,}1353+0{,}3679+1=1{,}503$; $\ln 1{,}503=0{,}4076$.
  \item $g(x)=1002+0{,}4076=1002{,}408$.
  \item Directamente: $e^{1000}\approx 10^{434}$ supera el máximo de doble precisión ($\approx 1{,}8\cdot 10^{308}$); se obtiene $+\infty$ y $g=+\infty$, un resultado falso.
\end{enumerate}
\resultado \textbf{$g(x)\approx 1002{,}408$}; el cálculo directo desborda. Además, se obtiene la softmax sin coste adicional:
\[\softmax(x)=\frac{(0{,}1353;\,0{,}3679;\,1)}{1{,}503}=(0{,}09003;\,0{,}2447;\,0{,}6652).\]
\comprobacion por la \cref{prop:lse}, $1002\le 1002{,}408\le 1002+\ln 3=1003{,}099$.
\end{solucion}

% ---------------------------------------------------------------------
\section{Tema 5}

\begin{solucion}{eje:5-ppl1}
\planteamiento propiedad (ii) de la \cref{prop:ppl}.
\pasos
\begin{enumerate}[label=\arabic*.]
  \item $\PPL=1$ equivale a $\mathcal J=0$.
  \item $\mathcal J$ es la media de los términos $-\ln p_t\ge 0$; es nula si y solo si todos lo son, es decir, $p_t=1$ para todo $t$.
\end{enumerate}
\resultado \textbf{El modelo predijo cada token del texto con probabilidad~1.}
\end{solucion}

\begin{solucion}{eje:5-perdida}
\datos $p=(1/2,\,1/4,\,1/4,\,1/8)$; $N=4$.
\planteamiento en base~2 los logaritmos son enteros; después se convierte a nats multiplicando por $\ln 2$.
\pasos
\begin{enumerate}[label=\arabic*.]
  \item $-\log_2 p_t=1,\,2,\,2,\,3$; suma $8$; media $\mathcal J=2$ bits.
  \item En nats: $2\ln 2=2\cdot 0{,}6931=1{,}386$.
  \item $\PPL=2^{2}=4$; equivalentemente, $e^{1{,}386}=4{,}000$.
\end{enumerate}
\resultado \textbf{$\mathcal J=2$ bits $=1{,}386$ nats; $\PPL=4$.}
\comprobacion por la \cref{prop:ppl}(i), $\PPL=(2\cdot 4\cdot 4\cdot 8)^{1/4}=256^{1/4}=4$.
\end{solucion}

\begin{solucion}{eje:5-digitos}
\datos $P(0)=0{,}82$; $P(d)=0{,}02$ para $d=1,\dots,9$; prueba con nueve ceros y un~3, $N=10$.
\planteamiento $\PPL=\exp\bigl(-\frac1N\sum_t\ln p_t\bigr)$.
\pasos
\begin{enumerate}[label=\arabic*.]
  \item $\ln 0{,}82=-0{,}1985$; para nueve ceros, $9\cdot(-0{,}1985)=-1{,}786$.
  \item $\ln 0{,}02=-3{,}912$.
  \item Suma $-1{,}786-3{,}912=-5{,}698$; pérdida media $5{,}698/10=0{,}5698$ nats.
  \item $\PPL=e^{0{,}5698}=1{,}768$. El modelo uniforme tiene $\PPL=10$ por la \cref{prop:ppl}(iii).
\end{enumerate}
\resultado \textbf{$\PPL\approx 1{,}768$ frente a~10 del modelo uniforme.}
\comprobacion $(0{,}82^{9}\cdot 0{,}02)^{-1/10}=(0{,}1676\cdot 0{,}02)^{-0{,}1}=0{,}003352^{-0{,}1}\approx 1{,}768$.
\end{solucion}

\begin{solucion}{eje:5-ppl}
\datos $p_1,\dots,p_N\in(0,1]$.
\planteamiento propiedades del logaritmo: $\ln(1/x)=-\ln x$, $\sum\ln a_t=\ln\prod a_t$, $a\ln x=\ln x^{a}$, $e^{\ln y}=y$.
\pasos
\begin{enumerate}[label=\arabic*.]
  \item $\mathcal J=\frac1N\sum_t\ln(1/p_t)=\frac1N\ln\prod_t(1/p_t)=\ln\bigl(\prod_t 1/p_t\bigr)^{1/N}$.
  \item $\PPL=e^{\mathcal J}=\bigl(\prod_t 1/p_t\bigr)^{1/N}$.
  \item Como $0<p_t\le 1$, cada $1/p_t\ge 1$; el producto es $\ge 1$ y, al ser $x\mapsto x^{1/N}$ creciente con $1^{1/N}=1$, también su raíz. La igualdad exige que todos los factores valgan~1.
  \item Si $p_t=1/K$ para todo $t$: $\PPL=(K^{N})^{1/N}=K$.
\end{enumerate}
\resultado \textbf{Las tres afirmaciones de la \cref{prop:ppl}.} $\blacksquare$
\comprobacion \cref{eje:5-perdida}: $(2\cdot 4\cdot 4\cdot 8)^{1/4}=4=2^{2}$.
\end{solucion}

\begin{solucion}{eje:5-lagrange}
\datos $p$ fija con $p(x)>0$ y $\sum_x p(x)=1$; se busca $q$ con $\sum_x q(x)=1$ que minimice $H(p,q)$.
\planteamiento minimización con una restricción de igualdad.
\pasos
\begin{enumerate}[label=\arabic*.]
  \item Lagrangiana: $\Lambda(q,\lambda)=-\sum_x p(x)\ln q(x)+\lambda\bigl(\sum_x q(x)-1\bigr)$.
  \item $\partial\Lambda/\partial q(x)=-p(x)/q(x)+\lambda=0$, luego $q(x)=p(x)/\lambda$.
  \item Restricción: $\sum_x p(x)/\lambda=1$, luego $\lambda=1$ y $q=p$.
  \item Es un mínimo global: $\partial^{2}H/\partial q(x)^{2}=p(x)/q(x)^{2}>0$ y las derivadas cruzadas son nulas, de modo que $H(p,\cdot)$ es estrictamente convexa sobre el conjunto convexo de distribuciones.
  \item Valor mínimo: $H(p,p)=H(p)$.
\end{enumerate}
\resultado \textbf{$H(p,q)$ alcanza su mínimo, $H(p)$, únicamente en $q=p$.} $\blacksquare$
\comprobacion con $p=(0{,}5;\,0{,}3;\,0{,}2)$: $H(p)=0{,}3466+0{,}3612+0{,}3219=1{,}030$ nats; con $q=(0{,}6;\,0{,}2;\,0{,}2)$ se obtiene $1{,}060$, y con la uniforme, $\ln 3=1{,}099$.
\end{solucion}
